% Chapter 6 --- Conclusions and Future Work
% Drafted to match the voice and evidence of Chapters 3-5.

\section{Summary of the Work}

This dissertation asked whether standard reinforcement learning algorithms
can learn profitable, genuinely state-dependent trading policies for a
single liquid asset. To answer this, a complete experimental system was
built: a daily trading environment for the SPY exchange-traded fund with
monthly episodes, transaction fees and action masking; a nine-feature state
representation with an optional drawdown extension; two reward
formulations; and three learning algorithms --- REINFORCE, Deep Q-learning
(DQN) and Proximal Policy Optimisation (PPO). Every component was covered
by an automated verification suite, and every design decision was made on
validation data before the 24 held-out test months were evaluated once.

The experiments followed a deliberate order. The fixed trade size was
calibrated first, so that the agents could reach a fair fraction of the
benchmark's market exposure. A behavioural diagnostic then established
which learned policies actually used their observations. Only after these
foundations were fixed were the agents compared against the benchmarks on
unseen data, followed by sensitivity experiments on the state
representation, the reward formulation and the level of transaction costs.

\section{Conclusions Against the Research Questions}

Chapter 1 posed four research questions. Each can now be answered with the
evidence of Chapter 5.

\textbf{Performance: no learned agent outperformed buy-and-hold.}
Buy-and-hold earned +\$180.25 per month on the held-out months. The best
learned agent, REINFORCE, earned +\$171.51 --- \$8.74 short --- while PPO
earned +\$127.35 and DQN +\$74.68. The gap is not explained by transaction
costs or by the action model, both of which were controlled. The test
period rose in 17 of 24 months, and staying invested was simply very hard
to beat. The honest conclusion is not that reinforcement learning failed,
but that the experiments found no evidence of learned skill beyond what an
unconditional buyer achieves.

\textbf{Behaviour: the algorithms split cleanly on whether they used the
state.} All six DQN seeds changed their actions in response to the observed
state. None of the twelve policy-gradient seeds (REINFORCE and PPO) did.
This behavioural finding reframes the performance ranking: REINFORCE's
near-benchmark return came from policies that converged to unconditional
buying, not from reading the market. Return alone would have hidden this
distinction entirely, which is the strongest argument this dissertation
makes for behavioural evaluation of learned traders.

\textbf{State representation: adding the drawdown feature changed little.}
Under the wealth-change reward, the ten-feature state moved REINFORCE by
$-\$4.12$ and DQN by $+\$3.13$ per month --- both far smaller than the seed
spreads of those methods --- and left the state-dependence results
unchanged. PPO's mean improved by \$32.62, but with a seed spread that
shrank from \$179.80 to \$94.00, the change cannot be attributed to the
feature with confidence. Information in the state is only useful if the
learning objective gives the agent a reason to use it.

\textbf{Reward: the risk-aware reward reduced risk by reducing
participation.} With the selected penalty weight of $\lambda = 0.10$, DQN's
mean drawdown fell 41\% (from 0.0140 to 0.0082), but its return fell 32\%
(from +\$77.81 to +\$52.90) and its exposure halved. Its Sharpe ratio fell
from 0.505 to 0.416, so risk-adjusted performance worsened. The penalty
taught the agents to hold less of the asset, not to time the market's
declines --- in March 2025, the month where stepping aside mattered most,
the state-dependent agent held \emph{more} than its average exposure and
took 46\% of the benchmark's loss.

\section{Critical Appraisal of the Project}

The most consequential decisions in this project were methodological rather
than algorithmic, and most of them were forced by problems found along the
way.

The verification-first approach paid for itself. The automated test suite
caught an implementation error in the buy-and-hold benchmark that
understated it by \$38.95 per month (21.6\%) --- an error that, left in
place, would have flipped the headline conclusion of this dissertation by
making the learned agents appear to beat the benchmark. All 21 verification
tests passed before any reported experiment was run. Similarly, the
trade-size calibration prevented a subtler mistake: under the original
$0.10\,C_0$ trade size, no agent could have reached more than 0.690
exposure against the benchmark's 0.910, and poor agent performance would
have been built into the environment rather than learned.

The behavioural diagnostic proved to be the most valuable analytical tool
in the study. Without it, REINFORCE's +\$171.51 would have read as a
near-success for reinforcement learning. With it, the same number reads
correctly: as the return of an unconditional buyer in a rising market.

Some parts of the project were harder than planned. Seed variance dominated
several comparisons --- PPO's spread of \$179.80 exceeded its own mean
return --- which meant that no single training run could support any
conclusion, and every configuration had to be trained six times from
scratch. The risk-aware reward also required more care than expected: the
penalty weight sweep showed a sharp behavioural cliff between $\lambda =
0.10$ and $0.25$, where exposure collapses from 0.344 to 0.074, so the
selection rule had to be designed around retaining participation rather
than around drawdown alone.

\section{Limitations}

The findings should be read within the following limits, stated before any
future work is proposed.

\begin{itemize}
\item
  \textbf{One asset, one regime.} All results concern SPY, and the
  held-out period (2024--2025) was generally rising. A market with
  prolonged declines would test the agents --- and the buy-and-hold
  benchmark --- very differently, and the conclusions here cannot be
  assumed to transfer.
\item
  \textbf{A fixed, discrete action model.} The agents trade fixed
  fractions of initial capital. The calibration in Section 5.3 made this
  fair against the benchmark, but a policy that can size its positions
  continuously might behave qualitatively differently.
\item
  \textbf{Six seeds, no formal significance testing.} Seed spreads are
  reported throughout and are often decisive, but the study stops short
  of hypothesis tests. With six seeds per configuration, small
  differences --- such as DQN's $+\$3.13$ state-representation change ---
  cannot be distinguished from noise, and are not claimed to be real.
\item
  \textbf{The penalty weight was selected using one algorithm.} The
  risk-aware reward's $\lambda$ was chosen with DQN, the only method
  whose policies respond to the state. The selected value may not be the
  best choice for the policy-gradient methods.
\item
  \textbf{The risk penalty targets only drawdown.} An agent can satisfy
  it by simply not participating. The results of Section 5.7 show this
  is what happened; a reward that relates return to risk, rather than
  penalising risk alone, was not tested.
\end{itemize}

\section{Future Work}

Each limitation points at a concrete next step, and the behavioural tools
built here would transfer to all of them.

\begin{itemize}
\item
  \textbf{Risk-efficiency objectives.} Replace the drawdown penalty with
  an objective that rewards return per unit of risk --- a Sharpe-style or
  drawdown-relative formulation --- so that reducing participation stops
  being the cheapest way to satisfy the reward. Section 5.7's finding
  that the penalty mainly bought lower exposure makes this the most
  direct continuation.
\item
  \textbf{Predictive, uncertainty-aware state information.} The drawdown
  feature describes the agent's own past losses, not the market's
  future. A state enriched with forecast information and an estimate of
  its own uncertainty would give a state-dependent learner something
  genuinely predictive to condition on, and would directly test whether
  DQN's willingness to use the state can be turned into timing skill.
\item
  \textbf{Harder market regimes.} Repeat the held-out evaluation over
  periods containing sustained declines, and move to walk-forward
  retraining so the agents see regime changes during training. The
  present study shows what happens in a rising market; the economically
  interesting question is what happens in a falling one.
\item
  \textbf{Richer action models.} Allow continuous position sizing, so
  that the trade-size calibration of Section 5.3 becomes unnecessary and
  the agent can express conviction rather than repetition.
\item
  \textbf{More seeds and formal statistics.} Scale the seed count and
  add significance testing so that small effects --- currently
  indistinguishable from noise --- can be confirmed or dismissed.
\end{itemize}

\section{Closing Remarks}

This dissertation set out to determine whether standard reinforcement
learning could beat the simplest investment strategy on unseen data, and
found that it could not: the best learned agent finished \$8.74 per month
behind buy-and-hold, and the agents that came closest did so by ignoring
their observations. The value of the work lies in how firmly that answer
is established. The comparison was made fair before it was made, the
policies were evaluated behaviourally as well as financially, and every
negative finding survived sensitivity checks over states, rewards and
transaction costs. The framework, diagnostics and corrected benchmarks
built here form a solid foundation on which the more ambitious agents of
future work can be evaluated just as honestly.
